\documentclass[conference]{IEEEtran}
\IEEEoverridecommandlockouts
\usepackage{cite}
\usepackage{amsmath,amssymb,amsfonts}
\usepackage{algorithmic}
\usepackage{graphicx}
\usepackage{textcomp}
\usepackage{xcolor}
\def\BibTeX{{\rm B\kern-.05em{\sc i\kern-.025em b}\kern-.08em
    T\kern-.1667em\lower.7ex\hbox{E}\kern-.125emX}}

\usepackage{listings}

\bibliographystyle{plain}

% --- Solidity Language Definition ---
\lstdefinelanguage{Solidity}{
  morekeywords={
    pragma, solidity, contract, library, interface, struct, event, modifier, function, returns,
    if, else, for, while, do, break, continue, return, emit, revert, require, assert,
    mapping, memory, storage, calldata, public, private, internal, external, view, pure, payable,
    address, bool, string, bytes, uint, uint8, uint16, uint32, uint64, uint128, uint256,
    int, int8, int16, int32, int64, int128, int256,
    new, delete, this, super, msg, block, tx
  },
  sensitive=true,
  morecomment=[l]{//},      % Single-line comments
  morecomment=[s]{/*}{*/}, % Multi-line comments
  morestring=[b]",          % Strings in double quotes
  morestring=[b]'           % Strings in single quotes
}

\definecolor{codegreen}{rgb}{0,0.6,0}
\definecolor{codegray}{rgb}{0.5,0.5,0.5}
\definecolor{codepurple}{rgb}{0.58,0,0.82}
\definecolor{backcolour}{rgb}{0.95,0.95,0.92}
% --- Doom One Light Color Definitions ---
% (Add this to your preamble)
\definecolor{doomBack}{RGB}{250, 250, 250}      % #fafafa
\definecolor{doomComment}{RGB}{156, 156, 156}   % #9c9c9c
\definecolor{doomKeyword}{RGB}{198, 120, 221}   % #c678dd (Purple/Magenta)
\definecolor{doomString}{RGB}{152, 195, 121}    % #98c379 (Green)
\definecolor{doomNumber}{RGB}{209, 154, 102}    % #d19a66 (Orange)
\definecolor{doomDefault}{RGB}{77, 77, 77}      % #4d4d4d (Default Text)
% Setup section. -----------------------------------|


\lstdefinestyle{soliditystyle}{
  language=Solidity,
  basicstyle=\ttfamily\small,
  backgroundcolor=\color{doomBack}, 
  commentstyle=\color{doomComment},
  keywordstyle=\color{doomKeyword},
  numberstyle=\tiny\color{doomNumber},
  stringstyle=\color{doomString},
  basicstyle=\ttfamily\footnotesize\color{doomDefault},
  breakatwhitespace=false,          
  breaklines=true,                  
  captionpos=b,                     
  keepspaces=true,                  
  numbers=left,                     
  numbersep=8pt,                   
  showspaces=false,                
  showstringspaces=false,
  showtabs=false,                   
  tabsize=2
}



\begin{document}

\title{Security and Software Application Project Report}

\author{\IEEEauthorblockN{Luca Sforza}
	\IEEEauthorblockA{\textit{Sapienza Università di Roma} \\
		sforza.2050030@studenti.uniroma1.it}
}

\maketitle

\begin{abstract}
	This paper illustrates the work performed for the exam project of the Security and Software Application course at Sapienza University of Rome.
	The main objective concerns the analysis and correction of vulnerabilities in Solidity smart contracts provided in the project specifications.
	The Agile methodology was adopted for development, dividing the work into incremental iterations to ensure constant control over code quality and system security invariants.
\end{abstract}

\section{Introduction}

The goal of this project consists of the analysis, correction, and securing of the smart contracts provided in the exam track.
The work focuses on identifying vulnerabilities present in the initial source code, verified through the definition and checking of the invariants described in the specifications.

The analysis phase was conducted following a \textit{Model Driven} approach, based on UML (\textit{Unified Modeling Language}), a formalism particularly suitable for modeling complex systems like those based on Solidity.

\section{Refactoring and Code Cleaning}

The initially provided code presented numerous structural and stylistic criticalities.
In particular, attribute visibility was set to the default value (often unintentionally) and indentation did not follow any recognized formatting standard.

To systematically identify and resolve these issues, static analysis tools (linters) were used.

\subsection{Development Environment}

For this project, the \textbf{Foundry} development framework was used, which provides a suite of essential tools for smart contract security and testing:

\begin{itemize}
	\item \textbf{Forge}: Used to compile, test, and deploy contracts.
	\item \textbf{Anvil}: A local Ethereum node for rapid testing.
	\item \textbf{Cast}: A tool for interacting with the contract via the command line.
\end{itemize}

Thanks to this framework, it is possible to format the Solidity code and use a linter to identify further errors.

\subsection{Analysis of Logical Problems}

Beyond syntactic issues, the code presented significant logical vulnerabilities. These were identified thanks to the \textit{Model Driven} approach.

The use of UML class diagrams allows for a natural translation into Solidity code. Literature proposes several meta-models to facilitate this transition.

The UML class diagram defines system constraints in the form of invariants, which must remain true during the entire execution of transactions.

A critical example encountered concerns the management of the taxpayer's age. Initially, this was stored as an integer to be manually incremented every year.
Such an approach is prone to errors and misalignments and is not gas-efficient.

Instead, age would be calculable if we saved the date of birth. Here a problem arises because Ethereum blockchain blocks use the UNIX timestamp, so Zero is 1970; thus, it is not possible to represent Taxpayers over 65.

Therefore, the approach used is to save the date (day, month, year), and the age is calculated via an approximation for leap year management.

A block's timestamp is converted into a date using an external dependency, \texttt{BokkyPooBahsDateTimeLibrary}. This increases the TCB (Trusted Code Base), but now the age is calculable.

This is not the only attribute that became calculable; for example, \texttt{isMarried} was made into a function.

\subsection{Taxpayer Factory}

Another criticality of the provided code is the \textbf{isContract} method. It serves to verify if an input address is a smart contract or not,
but this implementation does not ensure that a given smart contract has implemented the Taxpayer interface methods.
Therefore, a malicious contract could have this attribute, lying about which functions it implements.
Formally, any smart contract could potentially bring the system state into an invalid configuration using reentrancy attacks.

This can be resolved by creating a state-managed \textbf{FactoryTaxpayer} class. It has an owner that allows for instantiating taxpayers.
It registers them in its internal mapping.

The reference to this contract is saved in the taxpayer instances; taxpayers can query the factory to verify that a certain contract is a reliable Taxpayer using the method.

This guarantees security assuming the Factory owner is not compromised.

\section{Implementation of Constraints: Exercise 1}

\subsection{Analysis}

\begin{figure}[htbp]
	\centerline{\includegraphics[width=0.50\textwidth]{design/analysis/iter1.png}}
	\caption{
		UML class model of Agile iteration 1.
	}
	\label{uml:agile1}
\end{figure}

The first exercise requires the implementation of the marital reciprocity constraint ("spouse").

Figure \ref{uml:agile1} depicts the UML modeling of the system.
It defines the system invariant, which can be mostly constrained via Solidity language syntax.
In addition to the constraints imposed by the class diagram, external constraints exist.

The defined methods must atomically maintain the invariant during system iteration. They modify the system state and
must not bring the state into an invalid configuration.
Therefore, pre-conditions must be added to the methods to allow execution.

UML describes the system divided into two parts.
The part for interfacing with lotteries (which for the first iteration of the Agile methodology was ignored for now).
And a part for interfacing with Taxpayer.

Formally, the following first-order logic formula must be verified:

\begin{equation}
	\label{V_spouse}
	\forall a \in \text{Taxpayer}, \forall b \in \text{Taxpayer}: \text{spouse}(a,b) \implies \text{spouse}(b,a)
\end{equation}

If a taxpayer $a$ is registered as the spouse of $b$, then $b$ must necessarily have $a$ registered as their spouse.

If Taxpayer $A$ executes the \texttt{marry} method towards $B$, the contract must also force the update
of $B$'s state towards $A$ within the same transaction. This scenario generates a potential infinite recursion
if $B$ attempted to call \texttt{marry} back on $A$.

\subsection{Implementation}

The implemented solution involves introducing a new method called by the \texttt{marry} function,
and the pre-condition is that the marrying taxpayer is the caller of the \texttt{acceptMarriage} function.

The external constraints described in the diagram are guaranteed by the implementation of the `marry` method, where it atomically
marries two Taxpayers, rather than trivially being a setter.

The correctness of this implementation was validated using \textbf{Echidna} for fuzzing.

\begin{figure}
	\centering
	\begin{lstlisting}[style=soliditystyle]
function check_spouse(uint256 index) internal view returns (bool) {
    Taxpayer t = taxpayers[index];
    if (t.isMarried()) {
        return Taxpayer(t.getSpouse()).getSpouse() == address(t);
    }
    return true;
}
\end{lstlisting}
	\caption{Solidity design of invariant \ref{V_spouse}. This function is internal, so it will not be called directly
		by \textbf{Echidna}, but there is also a \textbf{public} function that iterates for each Taxpayer
		instantiated by the test contract to verify that this function is true; otherwise, it fails.}
	\label{code:V_spouse}
\end{figure}

\begin{figure}
	\centering
	\begin{lstlisting}[style=soliditystyle]
function marry(address newSpouse) public nonReentrant {
        require(newSpouse != address(0));
        require(!isMarried());
        require(newSpouse != address(this));
        require(newSpouse is a Taxpayer ERC165);

        marriage.spouse = newSpouse;
        marriage.maxAllowance = getTaxAllowance() + newSpouse.getTaxAllowance();

        newSpouse.acceptMarriage(address(this));
        assert(newSpouse.getSpouse() == address(this));
    }

    function acceptMarriage(address newSpouse) external nonReentrant {
        require(newSpouse is a Taxpayer ERC165);
        require(msg.sender == newSpouse);
        require(!isMarried());

        marriage.spouse = newSpouse;
        marriage.maxAllowance = this.getTaxAllowance() + ITaxpayer(newSpouse).getTaxAllowance();
    }
\end{lstlisting}
	\caption{
		Solidity design of the \textbf{marry} method. It calls the \texttt{acceptMarriage} function on the spouse's address.
	}
	\label{code:marry}
\end{figure}

The implementation can be read in figures \ref{code:V_spouse}
\ref{code:marry}.

Below is the screenshot where I execute fuzz testing on both the original contract and my described implementation.

\begin{figure}[htbp]
	\centerline{\includegraphics[width=0.50\textwidth]{design/report/funziona.jpg}}
	\caption{
		Correct Echidna output.
	}
	\label{echidna:corretto_es1}
\end{figure}

\begin{figure}[htbp]
	\centerline{\includegraphics[width=0.50\textwidth]{design/report/non_funziona_marry.jpg}}
	\caption{
		Correct Echidna output.
	}
	\label{echidna:sbagliato_es1}
\end{figure}

\subsection{Echidna Configuration}

To verify that the invariant is respected, there are two possible ways to do it with Echidna.

Defining the invariants directly in the smart contract of interest (\textbf{Taxpayer.sol}); this works as long as the invariant
must be verified locally in a Taxpayer instance.
However, invariant \ref{V_spouse} concerns multiple Taxpayer instances.

By default, \textbf{Echidna} instantiates the contract only once and then performs fuzz testing on the reference contract; it
will not try to instantiate the contract multiple times to verify that the invariant remains true. The most \textbf{Echidna} would do
is try to marry itself.

Therefore, a different solution must be found.

\textbf{Echidna} has the ability to define its internal configurations inside a \text{yaml} file
; a flag that can be activated is \textbf{allContracts}. If defined as \textbf{true}, then \textbf{Echidna} will perform fuzz testing not only
on the reference contract, but also on all contracts that the reference contract has instantiated.

Therefore, another possible solution is to write a new contract, called \textbf{EchidnaTesting}, where its constructor
instantiates various \textbf{Taxpayer} class objects and adds them to an array of Taxpayers.

The invariants, therefore, must not be added to the \textbf{Taxpayer} contract, but to the new test contract.

For the purpose of this project, 8 Taxpayers were instantiated, 2 of whom are over 65.

Thus, \textbf{Echidna} will try to marry different combinations of these 8 Taxpayers trying to make the invariant false.

Reentrancy attacks are not possible since all Taxpayers were created by a trusted party.


\section{Exercise 2}

\subsection{Analysis}
\begin{figure}
	\centering
	\begin{lstlisting}[style=soliditystyle]
  function marry(address newSpouse) public {
        require(newSpouse != address(0));
        require(!isMarried());
        require(newSpouse != address(this));
        require(factory.isTaxpayer(newSpouse));

        marriage.spouse = newSpouse;
        marriage.maxAllowance = this.getTaxAllowance() + ITaxpayer(newSpouse).getTaxAllowance();

        ITaxpayer(newSpouse).acceptMarriage(address(this));
        assert(newSpouse.getSpouse() == address(this));
        assert(newSpouse.getMaxTaxAllowance() == marriage.maxAllowance);
    }

    function acceptMarriage(address newSpouse) external {
        require(doesContractImplementInterface(newSpouse, type(ITaxpayer).interfaceId));
        require(msg.sender == newSpouse, "Caller must be the spouse");
        require(!isMarried());

        marriage.spouse = newSpouse;
        marriage.maxAllowance = this.getTaxAllowance() + ITaxpayer(newSpouse).getTaxAllowance();
    }
\end{lstlisting}
	\caption{
		Solidity design of the \textbf{marry} method. It calls the \texttt{acceptMarriage} function on the spouse's address.
		Prerequisites are modeled with the Solidity "require" construct. If the predicate is false, then
		the transaction will be reverted.
		Post-conditions are \textbf{asserts} since they are not needed in production, unlike
		requires which are necessary for correct software execution.
		However, in the testing phase, it is useful to have them because we can detect if the system state
		is entering an invalid configuration.
	}
	\label{code:marry2}
\end{figure}
\begin{figure}
	\centering
	\begin{lstlisting}[style=soliditystyle]
function check_spouse(uint256 index) internal view returns (bool) {
    Taxpayer t = taxpayers[index];
    if (t.isMarried()) {
        return Taxpayer(t.getSpouse()).getSpouse() == address(t);
    }
    return true;
}
\end{lstlisting}
	\caption{Solidity design of invariant \ref{V_spouse}. This function is internal, so it will not be called directly
		by \textbf{Echidna}, but there is also a \textbf{public} function that iterates for each Taxpayer
		instantiated by the test contract to verify that this function is true; otherwise, it fails.}
	\label{code:V_marry2}
\end{figure}

\begin{figure}[htbp]
	\centerline{\includegraphics[width=0.50\textwidth]{design/analysis/iter2.png}}
	\caption{
		The figure shows the second iteration of the Agile methodology analysis phase. Everything that was true
		for the class diagram in figure \ref{uml:agile1} is also true for this diagram.
	}
	\label{uml:agile2}
\end{figure}

With the conclusion of exercise 1, the first iteration of the Agile methodology was also concluded; therefore, the UML was modified
to model the specifications of exercise 2.

With the new specifications, it must be possible to move the tax allowance from one spouse to another, but with the constraint
that the sum of the two must equal the sum of the two tax allowances at the time of marriage.

This introduces a necessity to save this information when the marriage occurs.

Therefore, as shown in figure \ref{uml:agile2},
the information of the common marriage taxAllowance was added.

It was recorded during the marriage, and therefore when exchanging one's own taxAllowance,
the common taxAllowance must not be exceeded.

The constraint can be expressed through first-order logic in the following way.

\begin{align*}
	 & \forall a,b \in Taxpayer : married(a,b)       \\
	 & \implies a.taxAllowanceo + b.taxAllowance =   \\
	 & \text{Tax allowance at the start of Marriage}
\end{align*}

\subsection{Implementation}

As shown in figure \ref{code:marry2},
the new implementation of the marry method allows saving the maximum tax allowance amount of an entire household.

The code correctness was verified by implementing the constraint with Echidna as shown in figure \ref{code:V_marry2}.


\begin{figure}[htbp]
	\centerline{\includegraphics[width=0.50\textwidth]{design/report/funziona.jpg}}
	\caption{
		Correct Echidna output.
	}
	\label{echidna:corretto_es2}
\end{figure}

\begin{figure}[htbp]
	\centerline{\includegraphics[width=0.50\textwidth]{design/report/non_funziona_taxallowance.jpg}}
	\caption{
		Incorrect Echidna output.
	}
	\label{echidna:sbagliato_es2}
\end{figure}



\section{Exercise 3}

\begin{figure}[htbp]
	\centerline{\includegraphics[width=0.50\textwidth]{design/analysis/iter3.png}}
	\caption{
		In the design phase of Agile iteration number 3, the behavior for the case where a Taxpayer exceeds 65 years of age was modeled.
		Since a Taxpayer's age is calculable, the Taxpayer must claim the taxAllowance increase via the \texttt{redeemTaxAllowance} method. It has the pre-condition
		of being at least 65 years old. For a Taxpayer, it is in their interest and responsibility to claim the taxAllowance. If the Taxpayer is married, then the \texttt{maxTaxAllowance} must be updated
		via the \texttt{redeemTaxAllowanceOfSpouse} method.
	}
	\label{uml:agile3}
\end{figure}

\begin{figure}[htbp]
	\centerline{\includegraphics[width=0.50\textwidth]{design/analysis/redeem.png}}
	\caption{
		When a Taxpayer is instantiated, the contract state is `notRedeemed`. If a taxpayer is at least 65 years old, they can execute
		the \texttt{redeemTaxAllowance} function and obtain more tax allowance. Once in the `redeemed` state, it is not possible to exit.
	}
	\label{uml:redeem}
\end{figure}

With the end of Echidna testing, the second iteration of the Agile cycle also ended.

For the third part, we start again from the Model Driven phase of this project.

Since age is now calculable, we cannot automatically increase an individual's taxAllowance
when they turn 65. Additionally, if the Taxpayer is married, the global marriage taxAllowance must be modified.

Therefore, a method called `redeemTaxAllowance` was added.

It is in the interest of a taxpayer who has turned 65 to call this method as soon as they turn 65 to
update their taxAllowance. This requires a small fee to be paid to the blockchain when turning 65, but it only needs to be done
once in a lifetime.

Therefore, we can view the operation as a finite state machine as shown in figure \ref{uml:redeem}. % TODO: image

\section{Implementation}

\begin{figure}
	\label{code:V_reedem}
	\centering
	\begin{lstlisting}[style=soliditystyle]
function echidna_reedem_check(uint256 index) internal returns (bool) {
        Taxpayer t = taxpayers[index];
        if (t.isReedem()) {
            return t.age() >= 65;
        }
    }
  \end{lstlisting}
\end{figure}

The added invariant verifies that if you have redeemed, then you are at least 65 years old, as shown in figure \ref{code:V_reedem}.

Thanks to the pre-conditions of the \texttt{redeemTaxAllowance} method, the constraint is respected.

Obviously, if a Taxpayer increases their tax allowance, the common tax allowance value with their spouse must be updated.

All cases were checked with another invariant, not inserted in the report because it is too large.

% TODO: add image
\begin{figure}[htbp]
	\centerline{\includegraphics[width=0.50\textwidth]{design/report/funziona.jpg}}
	\caption{
		Correct Echidna output.
	}
	\label{echidna:corretto_es3}
\end{figure}

\begin{figure}[htbp]
	\centerline{\includegraphics[width=0.50\textwidth]{design/report/non_funziona_reedem.jpg}}
	\caption{
		Incorrect Echidna output.
	}
	\label{echidna:sbagliato_es3}
\end{figure}

\section{Exercise 4}

\subsection{Analysis}

\begin{figure}[htbp]
	\centerline{\includegraphics[width=0.50\textwidth]{design/analysis/iter4.png}}
	\caption{UML of the Analysis phase of Agile iteration number 4}
	\label{uml:agile4}
\end{figure}

\begin{figure}[htbp]
	\centerline{\includegraphics[width=0.50\textwidth]{design/analysis/lottery.png}}
	\caption{UML State Diagram of the Lottery. In the \texttt{notStarted} phase, no one can enter the lottery until it is activated. In the \texttt{started} state, Taxpayers
		can enter the lottery. When the finite state machine enters the \texttt{ending} state, the value of a future block is saved. If that block has arrived,
		then it is possible to draw the winner using the hash of the block where the transaction was made.
	}
	\label{uml:lottery}
\end{figure}

In the design phase of Agile iteration number 4, the Lottery modeling was added, as shown in figure \ref{uml:agile4}.
A lottery must implement the following methods: \texttt{startLottery}, \texttt{join}, \texttt{endLottery}, \texttt{selectWinner}.

The lottery also has an owner, which is the same as the FactoryTaxpayer.

Being the trusted party, they can decide when to end or start the lottery.

All references to time have been removed, since validators in the Ethereum blockchain could
manipulate that value; therefore, it must be the trusted party's responsibility to ensure they occur after a determined
period of time, using the real clock, not the blockchain's.

The implementation was improved so that the lottery is not "disposable," but instead can be
restarted as many times as desired; this is more gas-efficient compared to instantiating a new contract every time.

A finite state machine design pattern was implemented.
Each phase is represented by a state, where only specific operations can be performed in a certain state.

The lottery code was extensively modified because it had a very important limitation.

"Randomness" was managed by Taxpayers when they made their commits to the lottery, but user input cannot
be trusted.

The Ethereum blockchain was designed to solve this problem and have a much more reliable system for managing randomness.

When the lottery ends, a winner is not immediately declared; instead, the block number where
the transaction was inserted, added to a definable constant, is saved in an attribute.

Therefore, Taxpayers would have to convince a large portion of the Ethereum network to create blocks with a certain hash
to manipulate the lottery, which is unrealistic.


The behavior is shown in figure \ref{uml:lottery}.

This improves the lottery's information flow, making it more difficult for taxpayers to cheat.

\subsection{Implementation}

\begin{figure}
	\label{code:V_lottery2}
	\centering
	\begin{lstlisting}[style=soliditystyle]
function complete_iteration() public {
        create_lottery();
        joinAll();
        endLottery();
        wait(1 weeks); // NOTE: this make time pass one week
        selectWinner();
    }

    function echidna_check_lottery() public view returns (bool) {
        (address minAddress, uint256 minValue) = getMinWins();
        (address maxAddress, uint256 maxValue) = getMaxWins();

        if (wins[address(0)] > 0) return false;
        if (total_ends < 5) return true;
        return maxValue - minValue < 25;
    }
  \end{lstlisting}
	\caption{\texttt{getMinWins} gives the smallest win value of the distribution, while \texttt{getMaxWins} gives the largest. Thus constraining that after
		many turns, the variance must be within a certain range.}
\end{figure}


To constrain the lottery to function correctly, a constraint must be added
that the variance of games won by each participating taxpayer must be within a certain range.

The operation is shown in 17.% \ref{code:V_lottery2}.


A new Test contract was created, different from the previous one. Now the \texttt{allContract} flag is false.

Therefore, Echidna is forced to call only the public methods of the test contract.

The only public function is the one that allows Echidna to perform a full test round.


\begin{figure}[htbp]
	\centerline{\includegraphics[width=0.50\textwidth]{design/report/funziona_lottery.jpg}}
	\caption{
		Correct Echidna output.
	}
	\label{echidna:corretto_es4}
\end{figure}

\begin{figure}[htbp]
	\centerline{\includegraphics[width=0.50\textwidth]{design/report/non_funziona_lottery.jpg}}
	\caption{
		Incorrect Echidna output.
	}
	\label{echidna:sbagliato_es4}
\end{figure}


\bibliography{bibliography}

\end{document}
